\documentclass[12pt]{article}
\usepackage[margin=0.6in,top=0.85in,bottom=0.55in]{geometry}
\usepackage{lmodern}
\usepackage{microtype}
\usepackage{fancyhdr}
\usepackage{titlesec}
\usepackage{enumitem}
\usepackage{lastpage}
\usepackage[hidelinks]{hyperref}

\setlength{\parindent}{0pt}
\setlength{\parskip}{0.25em}
\setlength{\headheight}{26pt}
\titleformat{\section}{\large\bfseries\scshape}{}{0pt}{}
\pagestyle{fancy}
\fancyhf{}
\fancyhead[C]{\footnotesize Student Name: \hspace{2.3in} \quad Assignment ID: U1L4LE \quad Page \thepage\ of \pageref{LastPage}}
\renewcommand{\headrulewidth}{0.5pt}

\newcommand{\answerbox}[2]{%
\vspace{0.12em}
\noindent\textbf{#1}\par
\vspace{0.04em}
\noindent\fbox{%
\begin{minipage}[t][#2]{\dimexpr\textwidth-2\fboxsep-2\fboxrule}
\rule{\textwidth}{0pt}
\vspace{#2}
\end{minipage}}
\par\vspace{0.22em}}

\begin{document}

\begin{center}
{\LARGE\bfseries Week 4 Lit Example Worksheet}\par\vspace{0.05em}
{\large\scshape Julius Caesar: The Middle Term of Future Danger}\par\vspace{0.12em}
\rule{0.78\textwidth}{0.5pt}
\end{center}

\textbf{Purpose:} Apply Week 4's Whately method to Julius Caesar. Map practical syllogisms and test whether the middle term really connects the premises to the conclusion.

\textbf{Directions:} Use the Lit Example Reader. Quote briefly if useful. Your job is to map the reasoning, not retell the scene.

\section*{Part 1: Brutus's Main Argument}

\textbf{Question 1.1}\\[-0.15em]
In the serpent's egg passage, what conclusion is Brutus trying to justify? State it as a clear sentence.

\answerbox{ANSWER LABEL: Part 1 Q1}{2.0in}

\textbf{Question 1.2}\\[-0.15em]
Write one possible major premise and one possible minor premise behind Brutus's argument.

\answerbox{ANSWER LABEL: Part 1 Q2}{2.35in}

\newpage

\section*{Part 2: Terms in the Syllogism}

\textbf{Question 2.1}\\[-0.15em]
For Brutus's argument, name the minor term, major term, and middle term. Use ordinary language, not just S, P, and M.

\answerbox{ANSWER LABEL: Part 2 Q1}{2.35in}

\textbf{Question 2.2}\\[-0.15em]
Which premise is the pressure point: the major premise or the minor premise? Explain what Brutus would need to prove.

\answerbox{ANSWER LABEL: Part 2 Q2}{2.35in}

\newpage

\section*{Part 3: Cassius's Argument}

\textbf{Question 3.1}\\[-0.15em]
What conclusion does Cassius want Brutus to reach about Caesar or Rome? State one possible conclusion clearly.

\answerbox{ANSWER LABEL: Part 3 Q1}{2.0in}

\textbf{Question 3.2}\\[-0.15em]
Choose one possible middle term in Cassius's argument. Explain whether Cassius proves that Caesar belongs under that term.

\answerbox{ANSWER LABEL: Part 3 Q2}{2.65in}

\newpage

\section*{Part 4: Comparing the Arguments}

\textbf{Question 4.1}\\[-0.15em]
Which argument is more logically disciplined: Brutus's serpent's egg argument or Cassius's Colossus argument? Use the terms premise, conclusion, and middle term in your answer.

\answerbox{ANSWER LABEL: Part 4 Q1}{2.45in}

\textbf{Question 4.2}\\[-0.15em]
Put Whately's Week 4 question to the play: when a speaker argues from possible future danger, what middle term must be proved before violent prevention follows?

\answerbox{ANSWER LABEL: Part 4 Q2}{2.45in}

\end{document}
